\specialchapter{致读者}

\begin{enumerate}
\def\labelenumi{\arabic{enumi}.}
\item
  《数学原理》系列从数学的基础开始论述，并给出完整的证明。原则上，读者不需要具备特别的数学知识，只需熟悉数学推理，并具有一定的抽象思维能力。不过，本系列尤其面向至少掌握大学数学课程第一、二学年内容的读者。
\item
  我们选择的阐述方法是公理化的，通常从一般到特殊。证明的要求对主题材料施加严格固定的次序。因此，除非读者已经拥有相当广泛的数学知识，否则某些考察的用途不会立即显现。
\item
  本系列分为若干部，每部又分为若干章。法文版中已经全部或部分出版的各部列于下方。已有英文译本时，相应的英文书名列在括号中。在本卷中，凡有英文版可用，引用均指英文版，否则指法文版。
\end{enumerate}

集合论（集合论）记为 E（S）

代数（代数） --- A（A）

一般拓扑（一般拓扑） --- TG（GT）

（实变函数） --- FVR（FRV）

拓扑向量空间（拓扑向量空间） --- EVT（TVS）

积分（积分） --- INT（INT）

交换代数（交换代数） --- AC（CA）

微分流形与解析流形 --- VAR

（李群与李代数） --- LIE（LIE）

谱理论 --- TS

在前六部中（按上述顺序），正文中的每项陈述所假定的已知结果，仅限于同一章中已经讨论过的结果，或按下列顺序编排的前几章中的结果：S；A，第一至三章；GT，第一至三章；A，从第四章起；GT，从第四章起；FRV；TVS；INT。

从第七部起，读者通常会看到它与其他各部之间逻辑关系的明确说明（前六部始终假定为已知）。

\begin{enumerate}
\def\labelenumi{\arabic{enumi}.}
\setcounter{enumi}{3}
\item
  不过，我们有时在正文中插入一些例子，它们涉及读者可能已经知道、但本系列尚未讨论的事实。这些例子置于两个星号之间：*\ldots*。大多数读者无疑会发现，这些例子有助于理解正文。在其他情况下，置于*\ldots*之间的段落涉及本系列其他地方讨论的结果。我们希望读者能够验证其中没有循环论证。
\item
  每一章的逻辑框架由该章的定义、公理和定理构成。这些是后续使用时主要需要记住的部分。不太重要的结果，以及可以很容易由定理推出的结果，标作“命题”“引理”“推论”“评注”等。初读时可以省略的内容以小号字体印刷。对特别重要的定理，有时会以“附论”的名义给出评论。
\end{enumerate}

为了避免繁琐的重复，有时只在某一章或某章的某一节内有效的记号或缩写是方便的（例如，在只讨论交换环的章节中，“环”一词始终表示“交换环”）。这类约定总会明确说明，通常在它们出现的那一章开头。

\pandocbounded{\includegraphics[keepaspectratio]{images/part001_a52dd4e1.jpg}}

\begin{enumerate}
\def\labelenumi{\arabic{enumi}.}
\setcounter{enumi}{5}
\item
  正文中的某些段落旨在预先提醒读者注意严重错误。这些段落在页边以（“危险弯道”）的标志标出。
\item
  习题既用于使读者检验自己是否已消化正文，也用于提醒读者注意那些不适合放入正文、但仍然有趣的结果。最困难的习题带有符号 Ⅱ。
\item
  一般而言，我们遵循公认的术语，除非有充分理由偏离它们。
\item
  我们特别努力始终使用严格正确的语言，同时不牺牲简明性。我们尽可能在正文中指出语言滥用；没有这些指出，任何数学文本都可能陷入迂腐，更不用说难以阅读。
\item
  由于原则上正文由某个理论的教条式阐述构成，通常不包含对文献的引用。书目引用汇集在历史注记中。每篇历史注记之后的书目通常只包括对所讨论理论的发展最为重要的著作和原始论文，并不声称以任何方式达到完整。
\end{enumerate}

至于习题，我们通常认为标明其来源没有价值，因为它们取自许多不同来源（原始论文、教科书、习题集）。

\begin{enumerate}
\def\labelenumi{\arabic{enumi}.}
\setcounter{enumi}{10}
\tightlist
\item
  本系列某一部分的引用方式如下：
\end{enumerate}

\begin{enumerate}
\def\labelenumi{\alph{enumi})}
\item
  如果引用的是同一节（§）中陈述的定理、公理或定义，则以其编号引用。
\item
  如果它们出现在同一章的另一节中，则引用中还要标明该节。
\item
  如果它们出现在同一部的另一章中，则标明该章和该节。
\item
  如果它们出现在另一部中，则首先用该部书名的缩写引用该部。
\end{enumerate}

结果汇总以字母 R 引用；因此，S, R 表示“集合论的结果汇总”。

